\section{Conclusion}\label{sec:conclusion}
This paper formulated attention-aware task scheduling as a partially observable Markov decision process. In this formulation a worker's attention, fatigue and circadian state shape their capacity, rest is an explicit action, and the cognitive state is observed only through noisy estimates. We proposed a memory-augmented deep Q-network (AA-DQN) and evaluated it against nine rule-based schedulers, including grid-tuned EDF and SPT rules with periodic or state-triggered rest, using ten training seeds, 500 held-out workdays with common random numbers, and multiplicity-corrected tests.

Rules that ignore cognitive dynamics performed poorly, and every competitive policy rested in 15--20\% of the day. The pre-specified hybrid AA-DQN matched the return of the best tuned rule and achieved a higher on-time rate (\aaOnTimePct\% versus \bestHeurOnTimePct\%). A compact variant restricted to intensity actions, evaluated as a secondary analysis, exceeded the best rule in return (+\SecRetSPTThreshold), on-time rate (\AADQNThreeactOnTimePct\%) and difficulty-weighted on-time rate (\AADQNThreeactWOnTimePct\%). The learned policies are the only ones that combine high task throughput with protection of demanding tasks. They work intensively while the worker is fresh, schedule hard tasks when attention is high and rest around the post-lunch dip. Ablations showed that the value of learning grows sharply with the reliability of cognitive-state estimates, that memory is the critical architectural component, and that a moderate weight on fatigue improves both well-being and productivity.

The study is based on a behaviourally grounded but uncalibrated simulator. Its main limitation is therefore the gap to real workers, which future work will close through calibration with physiological and behavioural data, personalised policies, and user studies. The released simulator, agents and evaluation protocol provide a reproducible basis for that work and for fair comparisons of learned and hand-designed schedulers in human-centric systems.
